% !TEX root = ../metatime_monograph.tex
\chapter{The L-Constants}
\label{ch:l_constants}

Three geometric constants underpin the Metatime arithmetic:
\(L_3 = 7\), \(L_4 = 2\), and \(L_5 = 5\).  In this chapter we derive
each from the Collatz iteration on prime-indexed orbits and from the
twin-prime structure of the number 5.

\section{The Collatz Iteration}

The Collatz function (also known as the \(3n+1\) problem) is defined on
positive integers by:

\begin{equation}
f(n) = \begin{cases}
3n + 1 & \text{if } n \text{ is odd}, \\
n/2   & \text{if } n \text{ is even}.
\end{cases}
\label{eq:collatz}
\end{equation}

Starting from any positive integer, repeated application of \(f\) eventually
reaches the cycle \(4 \to 2 \to 1 \to 4\).  The Collatz depth of an integer
\(n\) is the number of steps required to reach 1.

\subsection{Collatz Depth of 3}

Starting from \(n = 3\):

\[
\begin{array}{c|c}
\text{Step} & \text{Value} \\ \hline
0 & 3 \text{ (odd)} \\
1 & 3\times 3 + 1 = 10 \text{ (even)} \\
2 & 10/2 = 5 \text{ (odd)} \\
3 & 3\times 5 + 1 = 16 \text{ (even)} \\
4 & 16/2 = 8 \text{ (even)} \\
5 & 8/2 = 4 \text{ (even)} \\
6 & 4/2 = 2 \text{ (even)} \\
7 & 2/2 = 1
\end{array}
\]

The trajectory is: \(3 \to 10 \to 5 \to 16 \to 8 \to 4 \to 2 \to 1\),
requiring 7 steps.  Hence:

\begin{equation}
L_3 \equiv \text{Collatz depth of 3} = 7.
\label{eq:L3}
\end{equation}

\section{The Twin-Prime Constants}

The twin-prime pair \((3,5)\) gives the remaining L-constants:

\begin{itemize}
\item The gap between the twin primes is \(5 - 3 = 2\), which defines the
annihilation constant:
\[
L_4 \equiv 2.
\]

\item The larger twin prime is 5, which defines the intention dimension:
\[
L_5 \equiv 5.
\]
\end{itemize}

Together these satisfy:

\[
L_4^{L_5} = 2^5 = 32,
\]

which is the dimension of the intention-phase space in the NOEMA tetrahedron
construction.

\section{Composite Constants}

The L-constants combine to form several composite quantities that appear
repeatedly in the sector formulas:

\[
\begin{array}{lcl}
L_3 + L_4 = 9, & L_3 - 1 = 6, & L_4/L_5 = 2/5 = 0.4, \\
L_3 + L_4 + L_5 = 14, & L_3^2 = 49, & L_4^2 = 4, \\
L_4^{L_5} = 32, & L_3 L_4^{L_5} = 224, & (L_4/L_3)^2 = 4/49.
\end{array}
\]

\section{Physical Interpretation}

The L-constants admit the following physical interpretations:

\begin{description}
\item[\(L_3 = 7\):] The structural depth of flavour space, corresponding to
the number of Collatz steps from 3 to 1.  Appears in denominators and as an
exponent in mixing angles (\(\sin\theta_{13} = 1/L_3\)).

\item[\(L_4 = 2\):] The annihilation constant, representing the twin-prime
gap.  Appears as the fundamental binary splitting factor.

\item[\(L_5 = 5\):] The intention dimension, the larger twin prime.
Appears in ratios with \(L_4\) to set mixing scales.
\end{description}

These constants together with \(\kappa\) and the quark primes generate every
numerical quantity in the Metatime framework.
